\documentclass[../../main]{subfiles}

\begin{document}

\section{Anéis}
\label{sec:matematica:algebra:estruturas-algebricas:aneis}

Nesta seção consideramos estruturas algébricas sobre a assinatura
\[
	\Sigma = \{+, \cdot : 2\},
\]
isto é, duas operações binárias interpretadas em um mesmo conjunto.

Uma estrutura dessa assinatura consiste em um conjunto não vazio $A$ munido de interpretações
\[
	+^A : A \times A \to A
	\quad \text{e} \quad
	\cdot^A : A \times A \to A.
\]

\subsection{Definição de anel}

Um anel é um modelo da assinatura $\Sigma = \{+, \cdot : 2\}$ que satisfaz os seguintes axiomas:

\subsubsection{Estrutura aditiva}

O par $(A,+)$ é um grupo abeliano, isto é:

\begin{itemize}
	\item (Associatividade)
	      \[
		      \forall a,b,c \in A,\quad (a+b)+c = a+(b+c);
	      \]

	\item (Elemento neutro)
	      \[
		      \exists 0 \in A,\quad \forall a \in A,\quad a+0 = 0+a = a;
	      \]

	\item (Inverso aditivo)
	      \[
		      \forall a \in A,\ \exists (-a) \in A,\quad a+(-a)=0;
	      \]

	\item (Comutatividade)
	      \[
		      \forall a,b \in A,\quad a+b = b+a.
	      \]
\end{itemize}

\subsubsection{Estrutura multiplicativa}

A multiplicação é associativa:
\[
	\forall a,b,c \in A,\quad (a \cdot b)\cdot c = a \cdot (b \cdot c).
\]

\subsubsection{Distributividade}

A multiplicação é distributiva em relação à adição:
\[
	\forall a,b,c \in A,
\]
\[
	a \cdot (b+c) = a \cdot b + a \cdot c,
\]
\[
	(a+b) \cdot c = a \cdot c + b \cdot c.
\]

\subsection{Anel com unidade}

Um anel é unitário se existe um elemento $1 \in A$ tal que:
\[
	\forall a \in A,\quad a \cdot 1 = 1 \cdot a = a.
\]

\subsection{Anel comutativo}

Um anel é comutativo se:
\[
	\forall a,b \in A,\quad a \cdot b = b \cdot a.
\]

\subsection{Observação estrutural}

Um anel pode ser visto como a interação de:

\begin{itemize}
	\item uma estrutura de grupo abeliano $(A,+)$;
	\item uma operação binária adicional $\cdot$;
	\item compatibilidade dada pela distributividade.
\end{itemize}

Essa visão será fundamental para a definição de corpos e para a construção de anéis quociente.

\end{document}
